% Generated draft: requires mathematical and source review.
\documentclass[11pt]{article}
\usepackage[margin=25mm]{geometry}
\usepackage{fontspec}
\setmainfont{Latin Modern Roman}
\usepackage{amsmath,amssymb,mathtools,unicode-math}
\setmathfont{Latin Modern Math}
\usepackage{enumitem,needspace,xurl,hyperref}
\setlength{\parindent}{0pt}
\setlength{\parskip}{5pt}
\setlength{\emergencystretch}{3em}
\newcommand{\N}{\mathbb{N}}
\newcommand{\Z}{\mathbb{Z}}
\newcommand{\Q}{\mathbb{Q}}
\newcommand{\R}{\mathbb{R}}
\newcommand{\C}{\mathbb{C}}
\newcommand{\F}{\mathbb{F}}
\newcommand{\A}{\mathbb{A}}
\newcommand{\PP}{\mathbb{P}}
\newcommand{\E}{\mathbb{E}}
\newcommand{\eps}{\varepsilon}
\newcommand{\dd}{\,\mathrm{d}}
\begin{document}
\section*{UnsolvedMath ID: 20003356}
\textbf{Source problem code:} catalogue entry
\subsection{Definitions and mathematical statement}
\textbf{Conjecture 54 (G. Margulis, A. Selberg), as extracted.} If \(\langle\Gamma^+,\Gamma^-\rangle\) is discrete, then it is an arithmetic lattice in \(G\). The excerpt omits the definitions of \(G\) and \(\Gamma^\pm\) and ends with the unfinished phrase \emph{It was observed in [151] that}. Thus this bare implication cannot be interpreted for arbitrary groups and subgroups.\par The attached aid reconstructs a higher-rank product setting. For an integer \(k\ge2\), take \(G=\mathrm{SL}_2(\R)^k\). For \(v=(v_1,\ldots,v_k)\in\R^k\), define\[u^+(v)=\left(\begin{pmatrix}1&v_i\\0&1\end{pmatrix}\right)_{i=1}^k,\qquad u^-(v)=\left(\begin{pmatrix}1&0\\v_i&1\end{pmatrix}\right)_{i=1}^k.\] Let \(L_+,L_-\subset\R^k\) be full parameter lattices, meaning discrete additive subgroups of rank \(k\) spanning \(\R^k\), and set \(\Gamma^\pm=u^\pm(L_\pm)\), \(\Delta=\langle\Gamma^+,\Gamma^-\rangle\). The aid additionally requires irreducibility, whose exact original formulation is absent. If expressed as density of the projection of each parameter lattice to every nonempty proper coordinate subproduct, that condition must be checked against the cited theorem; no claim here removes the irreducibility hypothesis or includes \(k=1\). An arithmetic lattice means a discrete finite-covolume subgroup commensurable, up to the relevant ambient conjugacy, with the indicated integral arithmetic group; irreducibility excludes product decompositions, and nonuniform means noncompact quotient.\par The aid attributes the following precise criterion, in its intended irreducible setting, to Benoist--Oh. Write \(x\odot y=(x_iy_i)_{i=1}^k\). The group \(\Delta\) is discrete if and only if there are a totally real number field \(K\) of degree \(k\) and an ordering \(\sigma=(\sigma_1,\ldots,\sigma_k)\) of all its real embeddings such that\[L_+\odot L_-\subset\sigma(\mathcal O_K),\quad\text{that is, }x\odot y\in\sigma(\mathcal O_K)\text{ for every }x\in L_+,\ y\in L_-.\] In that case the aid reports that \(\Delta\) is commensurable with a conjugate of \(\sigma(\mathrm{SL}_2(\mathcal O_K))\), and hence is a nonuniform irreducible arithmetic lattice in \(G\). It also reports that a later theorem of Benoist--Miquel independently applies. Precise theorem statements, bibliographic data, and the omitted hypothesis matching are not supplied, so these are attributed literature claims requiring verification.\par The aid's companion corollary says that, under the displayed field and integrality condition, any nonzero \(a\in L_+\) has all coordinates nonzero. Reciprocal coordinate rescaling gives \(a^{-1}\odot L_+\) and \(a\odot L_-\), both full \(\Z\)-lattices in the same field copy \(\sigma(K)\). Conjugation in \(\mathrm{GL}_2(\R)^k\) by \((\operatorname{diag}(a_i^{-1},1))_i\) realizes these reciprocal changes of upper and lower parameters; signs of \(a_i\) need not admit real square roots for an inner diagonal conjugation in \(G\). For fixed \(K,\sigma\) and any \(\Z\)-bases \(v_1,\ldots,v_k\) of \(L_+\) and \(w_1,\ldots,w_k\) of \(L_-\), the integrality condition is equivalent to the finite set of conditions \(v_i\odot w_j\in\sigma(\mathcal O_K)\) for all \(1\le i,j\le k\). This finite check presupposes a field and its embeddings; it is not an algorithm that finds them from arbitrary lattices.
\subsection{Short English statement}
The truncated conjecture predicts arithmeticity from discreteness of a group generated by opposite unipotent lattices. The attached aid reports its solution for irreducible parameter lattices in \(\mathrm{SL}_2(\R)^k\), \(k\ge2\), through a totally real field criterion. The original hypotheses and exact cited theorems remain to be checked.
\subsection{Sources}
\begin{itemize}\item Local UnsolvedMath registry, permanent ID 20003356, code catalogue entry; canonical \texttt{aim-other-notes.json}, \texttt{notes[162]}, conjecture tag.\item American Institute of Mathematics Problem Lists, Other domain, \emph{Emerging applications of measure rigidity}, source item 54, Conjecture 54 attributed to G. Margulis and A. Selberg; no section supplied. \url{https://aimath.org/WWN/measrigid/measrigid.pdf}.\item Source reference [151], whose observation and bibliographic entry are both missing after the truncated lead-in. The attached aid additionally names Benoist--Oh and Benoist--Miquel without supplying titles, theorem numbers, publication details, or URLs.\item Attached research attempt 1, \emph{Arithmeticity of groups generated by opposite unipotent lattices}, recorded 2026-08-14: result \texttt{full\_solution}, problem-at-run \texttt{solved}, rigor \texttt{proved}, novelty confidence low; registry \texttt{solved}, classification \texttt{SOLVED-IN-LITERATURE}; statement status \texttt{reconstructed\_unverified}. This machine-generated aid is not peer reviewed and needs independent verification before citation.\end{itemize}
\end{document}
